\documentclass[a4paper,10pt,twoside,openany]{article}

\usepackage[lang=hebrew]{maths}
\usepackage{hebrewdoc}
\usepackage{stylish}
\usepackage{lipsum}
\let\bs\blacksquare

\setlength{\parindent}{0pt}

%%%%%%%%%%%%
% Styling %
%%%%%%%%%%%%

\usepackage{enumitem}

%%%%%%%%%%%%%
% Counters  %
%%%%%%%%%%%%%

\setcounter{section}{1}     
            
%BIBLIOGRAPHY
\usepackage[
backend=biber,
style=alphabetic,
]{biblatex}
\addbibresource{bibliography.bib} %Imports bibliography file

%%%%%%%%%%
% Title  %
%%%%%%%%%%
\title{
אלגברה ב' - גיליון תרגילי בית 1 \\
מטריצות מייצגות, ערכים עצמיים, ולכסון
\\
\vspace{1cm}
\large{תאריך הגשה: 12.11.2026}
}
\date{}

\begin{document}
\maketitle

\begin{exercise}
יהי
$\mcal{E} = \prs{e_1, \ldots, e_n}$
הבסיס הסטנדרטי של
$\mathbb{F}^n$,
ונזכיר כי עבור בסיס
$\mcal{B}$
של מרחב וקטורי סוף־מימדי
$V$
מעל
$\mbb{F}$
קיים איזומורפיזם
\begin{align*}
\rho_\mcal{B} : V &\to \mbb{F}^n \\
\text{.} \hphantom{lalala} v &\mapsto \brs{v}_\mcal{B}
\end{align*}
הראו כי
$\rho_{\mcal{E}} : \mathbb{F}^n \to \mathbb{F}^n$
הינה העתקת הזהות.
\end{exercise}

\begin{exercise}
יהי
$\mcal{B} \coloneqq \prs{v_i}_{i \in [n]} = \prs{v_1, \ldots, v_n}$
בסיס למרחב וקטורי
$V$.
נתונה
$T \colon V \to V$
הפיכה המקיימת
\[\text{.} T\prs{v_1 + 2 v_2} = \sum_{i \in [n]} v_i\]
מצאו את סכום כל מקדמי המטריצה
$\brs{T^{-1}}_{\mcal{B}}$.

למשל, אם
$\brs{T^{-1}}_{\mcal{B}} = \pmat{1 & 2 \\ 3 & 4}$,
הסכום הוא
$1+2+3+4 = 10$.
\end{exercise}

\begin{exercise}
יהיו
$V,W$
מרחבים וקטוריים סוף־מימדיים מעל שדה
$\mbb{F}$
ותהי
$T \in \Hom_{\mbb{F}}\prs{V,W}$
העתקה לינארית חד־חד ערכית.

יהיו
\begin{align*}
\mcal{B} &= \prs{v_1, \ldots, v_n} \\
\mcal{C} &= \prs{u_1, \ldots, u_n}
\end{align*}
בסיסים של
$V$
ויהיו
\begin{align*}
\mcal{B}' &= \prs{T\prs{v_1}, \ldots, T\prs{v_n}} \\
\text{.} \mcal{C}' &= \prs{T\prs{u_1}, \ldots, T\prs{u_n}}
\end{align*}

\begin{enumerate}
\item
הראו כי
$\dim V = \dim \im T$.

\item הראו כי
$\mcal{B}', \mcal{C}'$
קבוצות בלתי תלויות לינאריות, והסיקו כי הן בסיסים של
$\im\prs{T} = \set{T\prs{v}}{v \in V}$.

\item
הראו כי
$M^\mcal{B}_\mcal{C} = M^{\mcal{B}'}_{\mcal{C}'}$.

\textbf{רמז:}
כיתבו את
$M^\mcal{B}_\mcal{C} e_i$
ואת
$M^{\mcal{B}'}_{\mcal{C}'} e_i$
כוקטורי קואורדינטות לפי הבסיסים
$\mcal{C}$
ו־%
$\mcal{C}'$
בהתאמה, והראו שמתקיים שוויון.
\end{enumerate}
\end{exercise}

\begin{exercise}
כדי ללכסן אופרטור
$T \in \End_\FF\prs{V}$
נסתכל לפעמים על בסיס
$\mcal{B}$
של
$V$
ונמצא מטריצה
$P$
הפיכה עבורה
$D \coloneqq P^{-1} \brs{T}_B P$
מטריצה אלכסונית.
כדי לקבל מכך בסיס
$\mcal{C}$
עבורו
$\brs{T}_{\mcal{C}}$
אלכסונית נרצה שיתקיים
$P = M^{\mcal{C}}_{\mcal{B}}$,
ואכן אז
\[D = P^{-1} \brs{T}_B P = \prs{M^{\mcal{C}}_{\mcal{B}}}^{-1} \brs{T}_B M^{\mcal{C}}_{\mcal{B}} = \brs{T}_{\mcal{C}}\]
לפי נוסחת מעבר בסיס.

\begin{enumerate}
\item יהי
$n \in \NN_+$,
תהי
$P \in \Mat_n\prs{\FF}$
מטריצה הפיכה ויהי
$\mcal{E} = \prs{e_1, \ldots, e_n}$
הבסיס הסטנדרטי של
$\FF^n$.
מיצאו בסיס
$\mcal{C}$
של
$\FF^n$
עבורו
$P = M^\mcal{C}_\mcal{E}$.

\item יהי
$V$
מרחב וקטורי ממימד
$n$
מעל
$\FF$
ויהי
$\mcal{B}$
בסיס של
$V$.
מיצאו בסיס
$\mcal{C}$
של
$V$
עבורו
$P = M^\mcal{C}_\mcal{B}$.

\textbf{רמז:}
היעזרו בתרגיל הקודם.

\item
יהי
$V = \RR_3\brs{x}$
המרחב הוקטורי של פולינומים ממעלה לכל היותר
$3$
שמקדמיהם ממשיים, יהי
\[\mcal{B} = \prs{1, 1-x, \prs{1-x}^2, \prs{1-x}^3}\]
ויהי
$T \in \End_\FF\prs{V}$
עבורו
\[\text{.} \brs{T}_\mcal{B} = \pmat{
1 & 0 & 1 & 0 \\
0 & 1 & 0 & 1 \\
1 & 0 & 1 & 0 \\
0 & 1 & 0 & 1
}\]
\begin{enumerate}
\item 
מיצאו מטריצות
$D,P \in \Mat_4\prs{\RR}$
כך ש־%
$D$
אלכסונית ו־%
$P$
הפיכה וכך שמתקיים
\[\text{.} P^{-1} \brs{T}_\mcal{B} P = D\]
\item
מיצאו בסיס
$\mcal{C}$
של
$V$
עבורו
$P = M^\mcal{C}_\mcal{B}$
והסיקו כי
$\brs{T}_\mcal{C} = D$.
\end{enumerate}

\end{enumerate}
\end{exercise}

\end{document}